\section{Discussion \& Conclusion}

We present a novel framework for 3D-aware aesthetic viewpoint suggestion that learns a 3D aesthetic field via feedforward 3D Gaussian Splatting, enabling geometry-grounded aesthetic reasoning and efficient aesthetic viewpoint discovery from sparse input views.
Despite its effectiveness, our framework still has room for improvement.
First, it relies on camera poses to build the aesthetic field. While such information can be obtained from COLMAP~\cite{schonberger2016structure} and is often available in robots/drones and smartphones~\cite{apple_arkit}, a pose-free variant would broaden the applicability of our method. This could be achieved by incorporating recent advances in pose-free methods~\cite{ye2024no,jiang2025anysplat}.
Second, the quality of the aesthetic field depends on the accuracy of the reconstructed geometry, which is affected by both the backbone’s ability and the input coverage. While the former can be improved by stronger geometry backbones~\cite{wang2025vggt}, the latter can be addressed by selecting views to ensure sufficient scene coverage~\cite{xiao2024nerf,chen2024gennbv}. 
Third, our viewpoint search is constrained to regions supported by the initial observations. A promising extension is an active-perception loop that acquires additional views in promising directions to expand the aesthetic field and enlarge the feasible search space.
% with minimal user effort, requiring only a single scene scan.